\section{Definitions: The Machines Inside the Words}
\chapterline

In Chapter~5, we treated mathematical claims as geometric territories: a universal statement promises that every object stays inside a safe region, and negating it means hunting for a single witness that steps outside. But this raises an essential question: what draws the boundary line in the first place? How do we determine whether a candidate number, function, or geometric shape belongs inside the territory or stands outside it?

In everyday language, definitions are largely descriptive. Dictionaries record how words are commonly used, capturing nuances, historical context, and loose associations. If someone asks what makes a cup a ``mug,'' one might mention a handle, a cylindrical ceramic shape, and its use for hot drinks. If a container lacks a handle or is made of glass, whether it counts as a mug remains open to informal debate. Everyday language tolerates imprecise boundaries.

In pure mathematics, a definition is not a descriptive summary; it is an exact specification. It establishes precise criteria that classify mathematical objects unambiguously. Every object either satisfies all stated conditions and belongs to the defined category, or fails at least one condition and does not.

\begin{figure}[htbp]
\centering
\begin{tikzpicture}[
    font=\small,
    machine/.style={draw=sourceblue!80, fill=sourceblue!5, rounded corners=6pt, thick, inner sep=10pt},
    gate/.style={draw=structurepurple!80, fill=white, rounded corners=3pt, thick, text width=2.4cm, align=center, inner sep=4pt},
    inout/.style={draw=gray!50, fill=white, rounded corners=3pt, text width=2.2cm, align=center, inner sep=4pt}
]

% Main machine boundary
\node[machine, minimum width=13.0cm, minimum height=3.6cm] (box) at (6.0, 0) {};
\node[anchor=north west, font=\footnotesize\scshape, text=sourceblue!90!black] at (-0.2, 1.6) {The Definition Machine: $\operatorname{Even}(n)$};

% Left: Candidate Inputs
\node[inout] (in1) at (1.0, 0.6) {\textbf{Candidate:} $n = 14$};
\node[inout] (in2) at (1.0, -0.6) {\textbf{Candidate:} $n = 9$};

% Gates inside machine
\node[gate] (g1) at (4.5, 0) {\textbf{\color{structurepurple}Gate 1: Domain}\\\scriptsize Is $n \in \Z$?};
\node[gate] (g2) at (7.8, 0) {\textbf{\color{structurepurple}Gate 2: Form}\\\scriptsize $\exists k \in \Z,$\\$n = 2k$};

% Output streams
\node[draw=targetgreen!80, fill=targetgreen!10, rounded corners=3pt, thick, text width=2.6cm, align=center, inner sep=5pt] (outT) at (11.0, 0.8) {
    \textbf{\color{targetgreen!90!black}ACCEPTED}\\\scriptsize Certified Even\\$14 = 2(7)$
};

\node[draw=trapred!70, fill=trapred!10, rounded corners=3pt, thick, text width=2.6cm, align=center, inner sep=5pt] (outF) at (11.0, -0.8) {
    \textbf{\color{trapred!90!black}REJECTED}\\\scriptsize Fails Gate 2\\$9 = 2(4) + 1$
};

% Connecting arrows
\draw[-{Latex[length=2.5mm]}, thick, gray!70] (in1.east) -- (g1.west);
\draw[-{Latex[length=2.5mm]}, thick, gray!70] (in2.east) -- (g1.west);

\draw[-{Latex[length=2.5mm]}, thick, structurepurple] (g1) -- node[above=2pt, font=\tiny\bfseries, text=structurepurple]{Yes} (g2);

\draw[-{Latex[length=2.5mm]}, thick, targetgreen!80!black] (g2.east) -- ++(0.3, 0) |- (outT.west);
\draw[-{Latex[length=2.5mm]}, thick, trapred!80!black] (g2.east) -- ++(0.3, 0) |- (outF.west);

\end{tikzpicture}
\caption{A mathematical definition viewed as a two-stage condition check. Candidate objects arrive from the left and must pass each requirement. Candidate $n = 14$ is an integer (Gate 1) and factors as $2 \times 7$ (Gate 2), so it satisfies both conditions and is classified as even. Candidate $n = 9$ is an integer but fails the factorization test, so it is not even.}
\label{06:fig-definition-machine-even}
\end{figure}

\begin{historybox}
Throughout the eighteenth century, calculus pioneers including Leonhard Euler and Joseph-Louis Lagrange treated definitions as descriptive summaries of geometric intuition. A \emph{continuous curve} was described informally as a line drawn across paper without lifting the quill, and mathematicians took for granted that any unbroken curve must possess a well-defined tangent line---an instantaneous slope or derivative---at almost every point. If you could sketch it without tears or leaps, smooth calculus tools were assumed to apply.

In 1872, Karl Weierstrass challenged this visual confidence by introducing a counterexample to the Berlin Academy: an infinite sum of cosine oscillations, adding infinitely many ripple waves of ever-smaller wavelengths and ever-steeper slopes. The resulting curve has no breaks or gaps, yet it is nowhere differentiable. At every single point, under every level of magnification, the curve consists entirely of sharp corners, never flattening into a line. Charles Hermite recoiled with ``fright and horror from this lamentable plague of functions without derivatives,'' while Henri Poincar\'e complained that logic had invented monsters to mock common sense.

Yet Bernard Bolzano (in 1817) and Weierstrass realized that the fault lay not in the functions, but in relying on human sketches. An ink line on paper has physical thickness that hides microscopic oscillations. Definitions could no longer be fuzzy summaries of visual intuition; they had to become rigid sorting machines governed by quantifiers ($\forall$ and $\exists$) and exact inequality contracts ($\eps$ and $\delta$). By defining continuity through arithmetic bounds rather than sketches, mathematics grounded proof in symbolic precision.

This shift remains vital today: in modern computer proof assistants such as Lean, Coq, and Isabelle, mathematical definitions are literally encoded as machine-checked sorting gates. A major frontier in modern mathematical reasoning is whether automated systems can learn to formulate the right minimal definitions on their own, rather than merely verifying human-crafted ones.
\end{historybox}

\subsection{The two operations: unpacking and packaging}

Every mathematical definition operates as a bidirectional pipeline. Depending on where you stand in a proof, you will use it in one of two modes:

\begin{enumerate}
\item \textbf{Unpacking (Using the Definition)}: You are \emph{given} that an object satisfies a definition. 
The definition allows you to replace the term with its explicit mathematical conditions:
\begin{quote}
``Since $n$ is even, there exists some integer $k$ such that $n = 2k$.''
\end{quote}
Here, the word ``even'' is replaced by a concrete algebraic equation ($n = 2k$) that you can substitute, manipulate, or factor.

\item \textbf{Packaging (Applying the Definition)}: You have performed algebraic manipulations and want to \emph{conclude} that an object satisfies a definition.
The definition specifies the exact entrance criteria you must satisfy:
\begin{quote}
``Since $m = 2(j^2 + j)$ where $j^2 + j \in \Z$, the number $m$ is of the form $2k$ for an integer $k$. Therefore, $m$ is even.''
\end{quote}
Here, you verify that the expression satisfies the conditions and conclude that $m$ is even.
\end{enumerate}

Early proof writing is largely the discipline of alternating cleanly between unpacking your hypotheses and packaging your conclusions, as illustrated in Figure~\ref{06:fig-unpack-package-pipeline}.

\begin{figure}[htbp]
\centering
\begin{tikzpicture}[
    font=\small,
    stage/.style={draw=gray!60, fill=white, rounded corners=5pt, thick, inner sep=6pt, text width=3.2cm, align=center},
    steparrow/.style={-{Latex[length=2.5mm]}, very thick}
]

% Stage 1: Unpack
\node[stage, draw=sourceblue!80, fill=sourceblue!5] (s1) at (1.6, 0) {
    \textbf{\color{sourceblue!90!black}1. Unpack Given}\\[3pt]
    \scriptsize \textbf{Sealed Box:} ``$n$ is even''\\[2pt]
    $\Downarrow$ \emph{open definition}\\[2pt]
    \textbf{Raw Algebra:}\\
    $\exists k \in \Z, \, n = 2k$
};

% Stage 2: Algebraic Manipulation
\node[stage, draw=structurepurple!80, fill=structurepurple!5, text width=3.6cm] (s2) at (6.8, 0) {
    \textbf{\color{structurepurple!90!black}2. Algebraic Engine}\\[3pt]
    \scriptsize Examine target quantity $n^2$:\\[2pt]
    $n^2 = (2k)^2 = 4k^2$\\[2pt]
    Factor out $2$:\\[2pt]
    $n^2 = 2 \cdot (2k^2)$
};

% Stage 3: Package
\node[stage, draw=targetgreen!80, fill=targetgreen!5] (s3) at (12.0, 0) {
    \textbf{\color{targetgreen!90!black}3. Package Goal}\\[3pt]
    \scriptsize Check Gate: Is $2k^2 \in \Z$? Yes!\\[2pt]
    $\Downarrow$ \emph{seal definition}\\[2pt]
    \textbf{Certified Label:}\\
    ``$n^2$ is even''
};

% Connecting Arrows
\draw[steparrow, sourceblue!80!black] (s1.east) -- node[above=2pt, font=\scriptsize\bfseries, text=sourceblue!90!black]{Unpack} (s2.west);
\draw[steparrow, targetgreen!80!black] (s2.east) -- node[above=2pt, font=\scriptsize\bfseries, text=targetgreen!90!black]{Package} (s3.west);

\end{tikzpicture}
\caption{Unpacking and packaging in a direct proof. On the left, we unpack the hypothesis: replacing the term ``even'' with the algebraic statement $n = 2k$ for some $k \in \Z$. In the center, we manipulate the expression algebraically into the form $2 \cdot (\text{integer})$. On the right, we verify that the condition is met and conclude that $n^2$ is even.}
\label{06:fig-unpack-package-pipeline}
\end{figure}

\subsection{Divisibility: a definition to handle with care}

Let us inspect a standard arithmetic definition that beginners often mangle:

\begin{definition}
Let $a, b \in \Z$. We say that $a$ \emph{divides} $b$, written $a \mid b$, if there exists an integer $k$ such that $b = ak$.
\end{definition}

In plain language, $a \mid b$ means that $b$ is an exact whole-number multiple of $a$. Imagine a measuring rod of length $a$: it can tile a segment of length $b$ perfectly, with no leftover remainder.

Notice the orientation of the symbols:
\begin{itemize}
\item $3 \mid 12$ is true because $12 = 3(4)$. Here $a = 3$, $b = 12$, and the witness integer is $k = 4$.
\item $12 \mid 3$ is false because there is no integer $k$ satisfying $3 = 12k$ (the only real solution is $k = 1/4$, which fails the integer gate requirement).
\end{itemize}

Figure~\ref{06:fig-divisibility-tiling} illustrates this measuring-rod model: $a$ divides $b$ precisely when rods of length $a$ tile the interval $[0, b]$ without gap or overlap.

\begin{figure}[htbp]
\centering
\begin{tikzpicture}[
    x=0.88cm, y=1cm,
    font=\small,
    rod/.style={draw=sourceblue!80, fill=sourceblue!15, rounded corners=2pt, thick, inner sep=0pt},
    failrod/.style={draw=trapred!80, fill=trapred!15, rounded corners=2pt, thick, inner sep=0pt},
    remrod/.style={draw=quotientorange!90, fill=quotientorange!20, dashed, rounded corners=2pt, thick, inner sep=0pt},
    axis/.style={-{Latex[length=2mm]}, thick, gray!60}
]

% Top case: 3 | 12
\node[anchor=west, font=\footnotesize\bfseries, text=sourceblue!90!black] at (0, 3.1) {Success: $3 \mid 12$ \quad ($a = 3$, $b = 12$)};
\draw[axis] (0, 2.0) -- (12.7, 2.0) node[right, font=\footnotesize, text=gray!70!black] {$x$};
\foreach \x in {0, 3, 6, 9, 12} {
    \draw[gray!60, thick] (\x, 1.9) -- (\x, 2.1);
    \node[below=3pt, font=\scriptsize, text=gray!80!black] at (\x, 1.9) {$\x$};
}
\draw[rod] (0, 2.2) rectangle ++(3, 0.5) node[pos=0.5, font=\scriptsize, text=sourceblue!90!black]{rod $a=3$};
\draw[rod] (3, 2.2) rectangle ++(3, 0.5) node[pos=0.5, font=\scriptsize, text=sourceblue!90!black]{rod $a=3$};
\draw[rod] (6, 2.2) rectangle ++(3, 0.5) node[pos=0.5, font=\scriptsize, text=sourceblue!90!black]{rod $a=3$};
\draw[rod] (9, 2.2) rectangle ++(3, 0.5) node[pos=0.5, font=\scriptsize, text=sourceblue!90!black]{rod $a=3$};
\node[anchor=west, font=\scriptsize, text=targetgreen!80!black] at (0, 1.15) {Exact tiling: $12 = 3 \cdot 4 \implies$ integer witness $k = 4 \in \Z$};

% Bottom case: 5 does not divide 12
\node[anchor=west, font=\footnotesize\bfseries, text=trapred!90!black] at (0, 0.3) {Failure: $5 \nmid 12$ \quad ($a = 5$, $b = 12$)};
\draw[axis] (0, -0.7) -- (12.7, -0.7) node[right, font=\footnotesize, text=gray!70!black] {$x$};
\foreach \x in {0, 5, 10, 12} {
    \draw[gray!60, thick] (\x, -0.8) -- (\x, -0.6);
    \node[below=3pt, font=\scriptsize, text=gray!80!black] at (\x, -0.8) {$\x$};
}
\draw[failrod] (0, -0.5) rectangle ++(5, 0.5) node[pos=0.5, font=\scriptsize, text=trapred!90!black]{rod $a=5$};
\draw[failrod] (5, -0.5) rectangle ++(5, 0.5) node[pos=0.5, font=\scriptsize, text=trapred!90!black]{rod $a=5$};
\draw[remrod] (10, -0.5) rectangle ++(2, 0.5) node[pos=0.5, font=\scriptsize, text=quotientorange!90!black]{gap $2$};
\node[anchor=west, font=\scriptsize, text=trapred!80!black] at (0, -1.6) {Incomplete tiling: $12 = 5 \cdot 2 + 2 \implies$ remainder $2$, no integer witness $k \in \Z$};

\end{tikzpicture}
\caption{The integer tiling interpretation of divisibility $a \mid b$. In the upper case, the interval $[0, 12]$ is tiled exactly by four rods of length $a = 3$, certifying that $3 \mid 12$ with integer witness $k = 4$. In the lower case, attempting to tile $[0, 12]$ with rods of length $a = 5$ leaves an incomplete gap of length $2$, so no integer multiplier exists and $5 \nmid 12$.}
\label{06:fig-divisibility-tiling}
\end{figure}

\begin{trapbox}
Do not confuse the vertical bar $a \mid b$ with the forward slash $a / b$.
\begin{itemize}
\item $a \mid b$ is a \textbf{statement} (an assertion that can be true or false).
\item $a / b$ is a \textbf{term} (an algebraic expression denoting a rational number).
\end{itemize}
Writing ``if $a / b$ then\ldots'' is a grammatical type error: you cannot place a numerical quantity into the condition slot of an if-then statement.
\end{trapbox}

Beginners often stumble on boundary cases involving zero and negative numbers because intuition confuses the relational statement $a \mid b$ with the algorithmic process of long division:
\begin{itemize}
\item \textbf{Does $3 \mid 0$?} Run it through the machine: does there exist an integer $k$ such that $0 = 3k$? Setting $k = 0$ works, since $0 = 3(0)$. Thus, $3 \mid 0$ is true! In fact, every integer $a$ divides $0$, because $0 = a \cdot 0$.
\item \textbf{Does $0 \mid 3$?} The machine asks: does there exist an integer $k$ such that $3 = 0 \cdot k$? Since $0 \cdot k = 0$ for every integer $k$, no such witness can ever exist. The machine firmly rejects it: $0 \nmid 3$.
\item \textbf{Negative divisors:} Does $-4 \mid 12$? Yes, because $12 = (-4)(-3)$, and $-3 \in \Z$. Signs are no obstacle for the divisibility machine.
\end{itemize}

\begin{recoverybox}
\textbf{Pause and think}: Before reading further, test these three candidate pairs mentally against the definition of divisibility:
(1) Does $6 \mid 24$? \quad (2) Does $24 \mid 6$? \quad (3) Does $0 \mid 0$?

\emph{Answers}: (1) True: $24 = 6(4)$. (2) False: $6 = 24k$ has no integer solution. (3) True under the formal definition: $0 = 0 \cdot k$ holds for any integer $k$ (for instance, $k = 1$). Note that while $0 \mid 0$ is a valid true relation, the fraction $0/0$ is undefined arithmetic.
\end{recoverybox}

\subsection{Why mathematicians demand minimal, non-redundant criteria}

Why are definitions designed the way they are? Consider how you might describe a square in everyday speech:
\begin{quote}
``A four-sided polygon with equal sides, four right angles, opposite sides parallel, and equal-length perpendicular diagonals that bisect each other.''
\end{quote}
That description is accurate, but redundant as a mathematical definition.

If every consequence were included in the definition of a square, proving that a shape is a square would require verifying every item on that list.

Mathematicians prefer \textbf{minimal, non-redundant definitions}:
\begin{definition}
A \emph{square} is a quadrilateral with four equal sides and one right angle.
\end{definition}
The remaining geometric properties---such as parallel opposite sides, three additional right angles, and perpendicular diagonals that bisect each other---are not part of the definition. They are \emph{theorems} deduced from this minimal set of conditions.

Keeping definitions minimal simplifies proofs:
\begin{itemize}
\item \textbf{Minimal Verification}: To prove an object qualifies, you need only verify the essential defining conditions.
\item \textbf{Deductive Utility}: Once established, all known theorems about that class of objects apply immediately.
\end{itemize}

\begin{trapbox}
\textbf{The Variable Collision Trap}: When a proof provides multiple hypotheses, you must open each box with a \emph{fresh variable name}. 

Suppose you are given that $a \mid b$ and $b \mid c$. A frequent beginner mistake is to write:
\begin{center}
``$b = ak$ and $c = bk$.'' \quad \textbf{(WRONG!)}
\end{center}
By reusing the letter $k$, you have accidentally claimed that the two multipliers are identical: $b/a = c/b$. For instance, $2 \mid 6$ (where the factor is $3$) and $6 \mid 24$ (where the factor is $4$). Reusing $k$ forces $3 = 4$, collapsing your proof into nonsense.

\textbf{Golden Rule}: Every time you unpack an existential definition ($\exists$), introduce a new, unused variable ($j, m, k_1, k_2$).
\end{trapbox}

\subsection{Worked proof: divisibility is transitive}

A relation is \emph{transitive} if it passes along chains: whenever $a$ relates to $b$, and $b$ relates to $c$, then $a$ relates directly to $c$ (like falling dominoes, or ancestral descent: if you are the child of your parent, and your parent is the child of your grandparent, you are the descendant of your grandparent).

Let us watch the unpacking and packaging pipeline in action on a complete, formal proof.

\begin{theorem}\label{06:thm-divisibility-transitive}
Let $a, b, c \in \Z$. If $a \mid b$ and $b \mid c$, then $a \mid c$.
\end{theorem}

\begin{proof}
Suppose $a \mid b$ and $b \mid c$. 

\emph{(Unpacking step)}: By the definition of divisibility applied to $a \mid b$, there exists an integer $j$ such that $b = aj$. Similarly, by definition applied to $b \mid c$, there exists an integer $m$ such that $c = bm$. Notice that we deliberately selected distinct variables $j$ and $m$ to avoid the Variable Collision Trap.

\emph{(Algebraic core)}: We want to show that $c$ is an integer multiple of $a$. Substituting the expression for $b$ into the equation for $c$, we have:
\[
c = bm = (aj)m = a(jm)
\]
by the associative property of integer multiplication. 

\emph{(Packaging step)}: Since $j$ and $m$ are integers, their product $k = jm$ is also an integer (the integers are closed under multiplication). Thus, we have expressed $c$ in the form $c = ak$ for some $k \in \Z$. By definition of divisibility, this means $a \mid c$.
\end{proof}

Every step is accounted for: we opened the definitions to reveal the algebra, manipulated the symbols lawfully, and closed the definition to reach our destination.

Figure~\ref{06:fig-transitivity-tiling} shows this substitution geometrically: the total segment $c$ is built from $m$ blocks of length $b$, and each block of length $b$ is composed of $j$ rods of length $a$.

\begin{figure}[htbp]
\centering
\begin{tikzpicture}[
    font=\small,
    cblock/.style={draw=structurepurple!80, fill=structurepurple!10, rounded corners=3pt, thick},
    bblock/.style={draw=targetgreen!80, fill=targetgreen!15, rounded corners=2pt, thick},
    ablock/.style={draw=sourceblue!80, fill=sourceblue!15, rounded corners=2pt, thick}
]

% Level 1: Entire quantity c
\draw[cblock] (0, 3.2) rectangle (10.8, 3.8);
\node[font=\footnotesize\bfseries, text=structurepurple!90!black] at (5.4, 3.5) {Quantity $c$ \quad ($c = bm$)};

% Level 2: Decomposed into m=3 blocks of size b
\draw[bblock] (0, 1.9) rectangle (3.55, 2.5) node[pos=0.5, font=\scriptsize, text=targetgreen!90!black] {Block 1: $b = aj$};
\draw[bblock] (3.625, 1.9) rectangle (7.175, 2.5) node[pos=0.5, font=\scriptsize, text=targetgreen!90!black] {Block 2: $b = aj$};
\draw[bblock] (7.25, 1.9) rectangle (10.8, 2.5) node[pos=0.5, font=\scriptsize, text=targetgreen!90!black] {Block 3: $b = aj$};
\node[anchor=west, font=\scriptsize\bfseries, text=targetgreen!80!black] at (11.0, 2.2) {$m = 3$ blocks};

% Level 3: Decomposed into j=2 rods of size a per block
\foreach \i in {0, 1, 2} {
    \draw[ablock] (\i*3.625, 0.6) rectangle (\i*3.625 + 1.72, 1.2) node[pos=0.5, font=\scriptsize, text=sourceblue!90!black] {rod $a$};
    \draw[ablock] (\i*3.625 + 1.83, 0.6) rectangle (\i*3.625 + 3.55, 1.2) node[pos=0.5, font=\scriptsize, text=sourceblue!90!black] {rod $a$};
}
\node[anchor=west, font=\scriptsize\bfseries, text=sourceblue!80!black] at (11.0, 0.9) {$jm = 6$ rods};

% Explanatory summary beneath
\node[anchor=north west, font=\scriptsize, text=gray!80!black, align=left] at (0, 0.1) {
    \textbf{Hypotheses:} $b = aj$ (with $j \in \Z$) \quad and \quad $c = bm$ (with $m \in \Z$)\\
    \textbf{Conclusion:} $c = (aj)m = a(jm)$, so $c$ is an integer multiple of $a$ $\implies a \mid c$.
};

\end{tikzpicture}
\caption{Geometric decomposition showing that divisibility is transitive ($a \mid b$ and $b \mid c \implies a \mid c$). The total quantity $c$ is partitioned into $m = 3$ blocks of size $b$. Because $a \mid b$, each individual block of size $b$ is itself partitioned into $j = 2$ rods of size $a$. Nesting these partitions tiles all of $c$ into $jm = 6$ rods of size $a$, certifying that $a \mid c$ with integer witness $k = jm$.}
\label{06:fig-transitivity-tiling}
\end{figure}

\subsection{From single machines to deductive chains}

Up to now, we have tested definitions one at a time. Most theorems, however, connect multiple definitions through conditional claims, such as: ``If an integer $n$ is a multiple of $4$, then $n$ is even.''

In Theorem~\ref{06:thm-divisibility-transitive}, the path from hypothesis to conclusion was direct. But consider the claim: ``If $n^2$ is even, then $n$ is even.'' Unpacking the hypothesis gives $n^2 = 2k$, which yields $n = \sqrt{2k}$---an expression that does not readily reveal the integer factors of $n$.

When a direct algebraic proof is difficult, alternative proof strategies become necessary. In Chapter~7 (\emph{If-Then Statements and Proof Strategy}), we examine the logical structure of conditional statements, exploring contrapositive proof and contradiction.

\subsection{Exercises}

\subsubsection*{Part A: Quick Mental Checks}
\begin{enumerate}
\item For each pair of integers $(a, b)$, decide whether $a \mid b$ is true or false:
\begin{enumerate}[label=(\alph*)]
\item $(7, 42)$
\item $(42, 7)$
\item $(-3, 15)$
\item $(6, 0)$
\item $(0, 6)$
\end{enumerate}

\item Identify the fatal flaw in this student's attempt to prove that the sum of two even integers is even:
\begin{quote}
``Let $x$ and $y$ be even integers. By definition of even, $x = 2k$ and $y = 2k$. Adding them gives $x + y = 2k + 2k = 4k = 2(2k)$. Since $2k$ is an integer, $x + y$ is of the form $2m$, so $x + y$ is even.''
\end{quote}

\item Determine whether each of the following is a \textbf{statement} (an assertion that can be true or false) or a \textbf{term} (an expression denoting a number or object):
\begin{enumerate}[label=(\alph*)]
\item $a \mid b$
\item $a / b$
\item $n = 2k$
\item $2k + 1$
\end{enumerate}
\end{enumerate}

\subsubsection*{Part B: Direct Unpacking and Packaging Practice}
\begin{enumerate}[resume]
\item Prove that the product of two odd integers is odd. Explicitly identify where you unpack the definition of odd and where you package it.
\item Prove that if $a \mid b$ and $a \mid c$, then $a \mid (b + c)$.
\item Prove that if $a \mid b$, then $a \mid (bx)$ for every integer $x$.
\item An integer $n$ is defined to be a \emph{multiple of $4$} if there exists an integer $k$ such that $n = 4k$. Prove that if $n$ is a multiple of $4$, then $n$ is even.
\item A real number $x$ is \emph{rational} ($x \in \Q$) if there exist integers $p, q \in \Z$ with $q \neq 0$ such that $x = p/q$. Prove that if $x, y \in \Q$, then $x + y \in \Q$.
\end{enumerate}

\subsubsection*{Part C: Deeper Exploration}
\begin{enumerate}[resume]
\item Suppose $a, b \in \Z$ are non-zero integers such that $a \mid b$ and $b \mid a$. Prove that $a = b$ or $a = -b$. Where in your argument does the fact that $a$ and $b$ are integers (rather than arbitrary real numbers) play an essential role?
\item An integer $p > 1$ is defined to be \emph{prime} if its only positive divisors are $1$ and $p$. Why do mathematicians deliberately demand $p > 1$, excluding the number $1$ from being prime? (Consider what would happen to the Fundamental Theorem of Arithmetic---which promises that every integer greater than $1$ has a \emph{unique} factorization into primes---if $1$ were included).
\end{enumerate}

\begin{walkbox}
On your next walk, notice how definitions in everyday life (such as ``organic food,'' ``planet,'' or ``full-time employee'') differ from mathematical definitions. When astronomers reclassified Pluto in 2006 by demanding that a planet must ``clear the neighborhood around its orbit,'' they replaced an informal description with explicit classification criteria. What were the specific conditions, and which condition did Pluto fail to satisfy?
\end{walkbox}

\newpage
